
The capability-LC alignment question is treated as a partial correlation problem. A multi-layer statistical analysis is designed to establish existence (P1), mechanism (P2), and population-level confirmation (P3).

\subsubsection{3.1 Dataset}

The publicly available AlpacaEval 2.0 leaderboard CSV (accessed August 2026) is used, containing pairwise preference results for 223 diverse LLMs compared against a GPT-4 Turbo reference on 805 instruction-following prompts. The three variables of interest are:

\begin{itemize}
\item \textbf{win\_rate}: Fraction of pairwise comparisons where human annotators prefer the model (range: approximately 3--80\%)
\item \textbf{length\_controlled\_winrate (LC\_winrate)}: Win rate after GLM-based length debiasing \citep{Dubois2024LengthControlled} (range: approximately 2--95\%)
\item \textbf{avg\_length}: Mean response token count per model (range: approximately 400--2200 tokens)
\end{itemize}

All 223 rows have complete data; no imputation was required.

\textbf{Verbosity separability check}: VIF = 1.764 for both win\_rate and avg\_length (threshold: VIF < 5.0), confirming that OLS coefficients are interpretable and that the low rank-correlation between win\_rate and avg\_length also supports the Spearman partial correlation analysis (Figure 4).

\subsubsection{3.2 Primary Analysis: Partial Correlation (H-E1)}

The primary quantity of interest is $\rho$(win\_rate, LC\_winrate | avg\_length) --- the Spearman partial correlation after controlling for avg\_length.

\textbf{Implementation}: \texttt{pingouin.partial\textbackslash \_corr(method='spearman', alternative='greater')} with covariate avg\_length. The \texttt{alternative='greater'} option implements a one-tailed test consistent with the pre-registered directional hypothesis $\rho$ > 0. Pre-registered threshold: |r\_partial| $\geq$ 0.15, p < 0.05. Bootstrap 95\% CI: 1000 resamples, random\_state = 42.

\subsubsection{3.3 Mechanism Analysis: Standardized OLS (H-M1)}

\textbf{Design}: LC\_winrate regressed on win\_rate\_std + avg\_length\_std with \texttt{sklearn.StandardScaler} preprocessing and \texttt{statsmodels.OLS}. Dominance criterion: |$\beta _win|$ > |$\beta _len|$. A verbosity-only baseline (LC\_winrate ~ avg\_length\_std) provides an $R^2$ comparison.

\subsubsection{3.4 FWL-Inspired Robustness Verification (H-M2)}

\textbf{Design}: avg\_length is regressed from both win\_rate and LC\_winrate using OLS; Spearman $\rho$ is computed on the residual pairs. The convergence criterion is delta = |$\rho_{Pingouin} - \rho_{resid}$| < 0.02.

Note on scope: The Frisch-Waugh-Lovell theorem guarantees exact algebraic equality between partial OLS coefficients and residualized OLS estimates for Pearson correlation. For Spearman partial correlation, residualized ranks provide an independent estimator that approximates this principle; the 1.1 percentage point convergence reported here is empirical evidence of robustness, not a theorem guarantee.

\subsubsection{3.5 Population-Level Confirmation: Kruskal-Wallis and Dunn Post-Hoc (H-M3, H-C1)}

\textbf{Design}: win\_rate quartiles (Q1--Q4; N = 56, 56, 55, 56). \texttt{scipy.stats.kruskal} applied to both LC\_winrate (H-C1) and $\Delta$ = LC\_winrate - win\_rate (H-M3). Post-hoc: \texttt{scikit\textbackslash \_posthocs.posthoc\textbackslash \_dunn} with Bonferroni correction. Effect size: $\epsilon ^2$ = (H - k + 1) / (N $-$ k). Bootstrap CI on Dunn Q1 vs Q4 p (1000 resamples).

Quartile boundaries: Q1 $\leq$ 5.19\% win\_rate; 5.19\% < Q2 $\leq$ 14.29\%; 14.29\% < Q3 $\leq$ 27.85\%; Q4 > 27.85\%.

\subsubsection{3.6 Sub-Hypothesis Structure}

\begin{table*}[htbp]
\centering
\resizebox{\dimexpr 2\columnwidth+\columnsep\relax}{!}{%
\begin{tabular}{llllllll}
\toprule
Hypothesis & Type & Gate & Test &  &  &  &  \\
\midrule
H-E1 & Existence & MUST\_WORK & Partial correlation > 0.15 &  &  &  &  \\
H-M1 & Mechanism & MUST\_WORK & \ & $\beta _win$\ & > \ & $\beta _len$\ & in standardized OLS \\
H-M2 & Mechanism & SHOULD\_WORK & FWL delta < 0.02 &  &  &  &  \\
H-M3 & Mechanism & SHOULD\_WORK & KW on $\Delta$ across quartiles, p < 0.05 &  &  &  &  \\
H-C1 & Condition & SHOULD\_WORK & KW on LC\_winrate AND Dunn Q1 vs Q4, both p < 0.05 &  &  &  &  \\
\bottomrule
\end{tabular}
}
\end{table*}

